\chapter{Trigonometry: The Unit Circle}\label{ch:g11:trigo}

Trigonometry begins in right triangles, but its true home is the
\emph{unit circle}, where \omterm{def:g11:trigo:cossin}{cosine} and \omterm{def:g11:trigo:cossin}{sine} become \omterm{def:g11:func:function}{functions} of an arbitrary
\omterm{def:g10:numbers:sets}{real number}. This chapter installs the \omterm{def:g11:trigo:radian}{radian}, the winding of the real
line around the circle, the exact values to know by heart, and the
symmetries that generate all the classical identities. The study of
$\cos$ and $\sin$ as \omterm{def:g11:func:function}{functions} --- variations, \omterm{def:g11:deriv:derivative}{derivatives} --- is carried
out in \cref{ch:g12:trigo}.

\section{Radians and the winding}

\begin{definition}[Radian]\label{def:g11:trigo:radian}
Let $\mathcal C$ be the circle of radius $1$ centered at the origin $O$,
and $I(1, 0)$. The \emph{radian}\index{radian} measure of an angle at the
center of $\mathcal C$ is the length of the arc it cuts on $\mathcal C$.
Since the full circle has circumference $2\pi$:
\[
360^\circ = 2\pi \text{ rad},
\qquad
180^\circ = \pi \text{ rad},
\qquad
1 \text{ rad} = \frac{180^\circ}{\pi} \approx 57.3^\circ .
\]
\end{definition}

\begin{method}[Converting]\label{met:g11:trigo:convert}
Degrees and \omterm{def:g11:trigo:radian}{radians} are proportional: multiply by $\frac{\pi}{180}$ to go
from degrees to \omterm{def:g11:trigo:radian}{radians}, by $\frac{180}{\pi}$ the other way. For instance
$60^\circ = 60 \times \frac{\pi}{180} = \frac{\pi}{3}$, and
$\frac{3\pi}{4} = \frac{3 \times 180^\circ}{4} = 135^\circ$.
\end{method}

\begin{definition}[Winding the line on the circle]\label{def:g11:trigo:winding}
To each \omterm{def:g10:numbers:sets}{real number} $t$, associate the point $M(t)$ of $\mathcal C$
obtained by walking a distance $\abs t$ along the circle from $I$,
counterclockwise if $t \geq 0$, clockwise if $t < 0$. Since the
circumference is $2\pi$, the numbers $t$ and $t + 2k\pi$ ($k \in \Z$)
reach the same point: $M(t + 2k\pi) = M(t)$.
\end{definition}

\begin{example}
$M(0) = I$; $M\!\left(\frac{\pi}{2}\right) = (0, 1)$, the top of the
circle; $M(\pi) = (-1, 0)$; $M\!\left(-\frac{\pi}{2}\right) = (0, -1)$;
and $M\!\left(\frac{9\pi}{4}\right) = M\!\left(\frac{\pi}{4} +
2\pi\right) = M\!\left(\frac{\pi}{4}\right)$.
\end{example}

\section{Cosine and sine}

\begin{definition}[Cosine and sine]\label{def:g11:trigo:cossin}
For every real $t$, the \emph{cosine}\index{cosine} and
\emph{sine}\index{sine} of $t$ are the \omterm{def:g10:coordgeom:system}{coordinates} of the point $M(t)$ of
the unit circle:
\[
M(t) = (\cos t,\ \sin t).
\]
\end{definition}

\begin{omfigure}
\begin{tikzpicture}[scale=2.0]
  \draw[->] (-1.3,0) -- (1.4,0) node[right] {$x$};
  \draw[->] (0,-1.3) -- (0,1.4) node[above] {$y$};
  \draw[gray] (0,0) circle (1);
  \coordinate (O) at (0,0);
  \coordinate (M) at (55:1);
  \draw[thick,omDef] (O) -- (M);
  \draw[dashed] (M) -- (M |- O) node[below, font=\small] {$\cos t$};
  \draw[dashed] (M) -- (M -| O) node[left, font=\small] {$\sin t$};
  \draw[->,omThm,thick] (0.35,0) arc (0:55:0.35);
  \node[omThm, font=\small] at (27:0.55) {$t$};
  \fill (M) circle (0.7pt) node[above right] {$M(t)$};
  \fill (1,0) circle (0.7pt) node[below right] {$I$};
\end{tikzpicture}

{\small The unit circle: the \omterm{def:g10:numbers:sets}{real number} $t$, wound counterclockwise from
$I$, lands at the point $M(t)$ whose \omterm{def:g10:coordgeom:system}{coordinates} \emph{define} $\cos t$
and $\sin t$.}
\end{omfigure}

\begin{proposition}[First properties]\label{prop:g11:trigo:first}
For all $t \in \R$ and $k \in \Z$:
\[
-1 \leq \cos t \leq 1, \qquad
-1 \leq \sin t \leq 1, \qquad
\cos^2 t + \sin^2 t = 1,
\]
\[
\cos(t + 2k\pi) = \cos t, \qquad \sin(t + 2k\pi) = \sin t .
\]
\end{proposition}

\begin{proof}
$M(t)$ lies on the unit circle, so its \omterm{def:g10:coordgeom:system}{coordinates} are between $-1$ and
$1$, and they satisfy the circle's \omterm{def:g10:algebra:equation}{equation} $x^2 + y^2 = 1$, which is the
identity $\cos^2 t + \sin^2 t = 1$. Periodicity restates
$M(t + 2k\pi) = M(t)$ (\cref{def:g11:trigo:winding}).
\end{proof}

\begin{example}\label{ex:g11:trigo:identity}
If $\sin t = \frac35$ and $t \in \intcc{\frac\pi2}{\pi}$ (second
quadrant), then $\cos^2 t = 1 - \frac{9}{25} = \frac{16}{25}$, so
$\cos t = \pm\frac45$; in the second quadrant the \omterm{def:g10:coordgeom:system}{abscissa} is negative,
hence $\cos t = -\frac45$.
\end{example}

The values to know by heart:
\begin{center}
\begin{tabular}{c|ccccc}
$t$ & $0$ & $\dfrac{\pi}{6}$ & $\dfrac{\pi}{4}$ & $\dfrac{\pi}{3}$ & $\dfrac{\pi}{2}$\\[6pt]
\hline\rule{0pt}{14pt}%
$\cos t$ & $1$ & $\dfrac{\sqrt3}{2}$ & $\dfrac{\sqrt2}{2}$ & $\dfrac12$ & $0$\\[6pt]
$\sin t$ & $0$ & $\dfrac12$ & $\dfrac{\sqrt2}{2}$ & $\dfrac{\sqrt3}{2}$ & $1$
\end{tabular}
\end{center}

\begin{remark}
A memory aid: the \omterm{def:g11:trigo:cossin}{sines} read
$\frac{\sqrt0}{2}, \frac{\sqrt1}{2}, \frac{\sqrt2}{2}, \frac{\sqrt3}{2},
\frac{\sqrt4}{2}$, and the \omterm{def:g11:trigo:cossin}{cosines} are the same list reversed.
\end{remark}

\section{Associated angles}

\begin{proposition}[Associated angles]\label{prop:g11:trigo:associated}
\index{associated angles}
For all $t \in \R$:
\[
\begin{aligned}
\cos(-t) &= \cos t, & \sin(-t) &= -\sin t,\\
\cos(\pi - t) &= -\cos t, & \sin(\pi - t) &= \sin t,\\
\cos(\pi + t) &= -\cos t, & \sin(\pi + t) &= -\sin t,\\
\cos\left(\tfrac{\pi}{2} - t\right) &= \sin t, &
\sin\left(\tfrac{\pi}{2} - t\right) &= \cos t .
\end{aligned}
\]
\end{proposition}

\begin{proof}
Each line expresses a symmetry of the circle.

$M(-t)$ is the reflection of $M(t)$ in the $x$-axis: the \omterm{def:g10:coordgeom:system}{abscissa} is
unchanged, the \omterm{def:g10:coordgeom:system}{ordinate} changes sign.

$M(\pi - t)$ is the reflection of $M(t)$ in the $y$-axis: walking
backwards from $\pi$ by $t$ lands opposite (horizontally) to walking
forwards from $0$ by $t$. The \omterm{def:g10:coordgeom:system}{abscissa} changes sign, the \omterm{def:g10:coordgeom:system}{ordinate} is
unchanged.

$M(\pi + t)$ is diametrically opposite $M(t)$ (half a turn): both
\omterm{def:g10:coordgeom:system}{coordinates} change sign.

$M\!\left(\frac\pi2 - t\right)$ is the reflection of $M(t)$ in the
diagonal line $y = x$, which swaps the two \omterm{def:g10:coordgeom:system}{coordinates}.
\end{proof}

\begin{omfigure}
\begin{tikzpicture}[scale=2.0]
  \draw[->] (-1.35,0) -- (1.45,0) node[right] {$x$};
  \draw[->] (0,-1.35) -- (0,1.45) node[above] {$y$};
  \draw[gray] (0,0) circle (1);
  \fill (40:1) circle (0.7pt) node[above right, font=\small] {$M(t)$};
  \fill (-40:1) circle (0.7pt) node[below right, font=\small] {$M(-t)$};
  \fill (140:1) circle (0.7pt) node[above left, font=\small]
    {$M(\pi - t)$};
  \fill (220:1) circle (0.7pt) node[below left, font=\small]
    {$M(\pi + t)$};
  \draw[black, dashed, thin] (40:1) -- (-40:1);
  \draw[black, dashed, thin] (40:1) -- (140:1);
  \draw[black, dashed, thin] (140:1) -- (220:1);
  \draw[omDef, thin] (0,0) -- (40:1);
  \draw[omDef, thin] (0,0) -- (140:1);
  \draw[omDef, thin] (0,0) -- (220:1);
  \draw[omDef, thin] (0,0) -- (-40:1);
\end{tikzpicture}

{\small The four associated points: $M(-t)$ mirrors $M(t)$ in the
$x$-axis, $M(\pi - t)$ in the $y$-axis, and $M(\pi + t)$ is diametrically
opposite. Every identity of \cref{prop:g11:trigo:associated} is read off
this picture.}
\end{omfigure}

\begin{example}
$\cos\frac{2\pi}{3} = \cos\left(\pi - \frac\pi3\right)
= -\cos\frac\pi3 = -\frac12$, and
$\sin\frac{7\pi}{6} = \sin\left(\pi + \frac\pi6\right)
= -\sin\frac\pi6 = -\frac12$. The table of five values, extended by the
symmetries, covers the whole circle.
\end{example}

\section{Solving trigonometric equations}

\begin{theorem}[The equations $\cos t = \cos a$ and $\sin t = \sin a$]
\label{thm:g11:trigo:equations}
For \omterm{def:g10:numbers:sets}{real numbers} $t$ and $a$:
\[
\cos t = \cos a \iff t = a + 2k\pi \ \text{or}\ t = -a + 2k\pi
\quad (k \in \Z),
\]
\[
\sin t = \sin a \iff t = a + 2k\pi \ \text{or}\ t = \pi - a + 2k\pi
\quad (k \in \Z).
\]
\end{theorem}

\begin{proof}
Two points of the unit circle have the same \omterm{def:g10:coordgeom:system}{abscissa} exactly when they are
equal or reflections of each other in the $x$-axis; by
\cref{prop:g11:trigo:associated}, the reflection of $M(a)$ is $M(-a)$. So
$\cos t = \cos a$ means $M(t) = M(a)$ or $M(t) = M(-a)$, i.e.\
$t = \pm a$ up to a multiple of $2\pi$. Similarly, equal \omterm{def:g10:coordgeom:system}{ordinates} mean
equal points or reflections in the $y$-axis, and the reflection of $M(a)$
is $M(\pi - a)$.
\end{proof}

\begin{method}[Solving $\cos t = c$ on an interval]\label{met:g11:trigo:solve}
\begin{enumerate}
  \item Find \emph{one} angle $a$ with $\cos a = c$, from the table of
        known values.
  \item Write the general solutions $t = \pm a + 2k\pi$
        (\cref{thm:g11:trigo:equations}).
  \item Choose the \omterm{def:g10:numbers:sets}{integers} $k$ that land inside the requested \omterm{def:g10:numbers:interval}{interval},
        and list the solutions.
\end{enumerate}
Proceed likewise for $\sin t = c$ with $t = a + 2k\pi$ or
$t = \pi - a + 2k\pi$.
\end{method}

\begin{example}\label{ex:g11:trigo:solve}
Solve $\cos t = \frac12$ on $\intoc{-\pi}{\pi}$. One solution of
$\cos a = \frac12$ is $a = \frac{\pi}{3}$. The general solutions are
$t = \frac\pi3 + 2k\pi$ and $t = -\frac\pi3 + 2k\pi$. Inside
$\intoc{-\pi}{\pi}$, only $k = 0$ contributes:
$t \in \left\{-\frac\pi3,\ \frac\pi3\right\}$.

Solve $\sin t = -\frac{\sqrt2}{2}$ on $\intco{0}{2\pi}$. One angle is
$a = -\frac\pi4$; the solutions are $t = -\frac\pi4 + 2k\pi$ and
$t = \pi + \frac\pi4 + 2k\pi = \frac{5\pi}{4} + 2k\pi$. In
$\intco{0}{2\pi}$: $t = \frac{7\pi}{4}$ (from $k = 1$) and
$t = \frac{5\pi}{4}$.
\end{example}

\begin{omfigure}
\begin{tikzpicture}[scale=2.0]
  \draw[->] (-1.35,0) -- (1.45,0) node[right] {$x$};
  \draw[->] (0,-1.2) -- (0,1.3) node[above] {$y$};
  \draw[gray] (0,0) circle (1);
  \draw[omProp, thick] (0.5,-1.05) -- (0.5,1.15)
    node[above, font=\small] {$x = \tfrac12$};
  \fill (60:1) circle (0.7pt)
    node[above right, font=\small] {$M\!\left(\tfrac\pi3\right)$};
  \fill (-60:1) circle (0.7pt)
    node[below right, font=\small] {$M\!\left(-\tfrac\pi3\right)$};
  \draw[omDef, thin] (0,0) -- (60:1);
  \draw[omDef, thin] (0,0) -- (-60:1);
\end{tikzpicture}

{\small Solving $\cos t = \frac12$: the vertical line $x = \frac12$
(orange) cuts the unit circle in two points, symmetric in the $x$-axis
--- hence the two solution families $t = \pm\frac\pi3 + 2k\pi$.}
\end{omfigure}

\section{Exercises}

\begin{exercise}[$\star$]\label{exo:g11:trigo:1}
Convert to \omterm{def:g11:trigo:radian}{radians}: $30^\circ$, $45^\circ$, $120^\circ$, $270^\circ$.
Convert to degrees: $\frac{\pi}{5}$, $\frac{5\pi}{6}$, $\frac{7\pi}{4}$,
$\frac{2\pi}{9}$.
\end{exercise}

\begin{exercise}[$\star$]\label{exo:g11:trigo:2}
Place on the unit circle the points $M(t)$ for
\[
t = \frac{2\pi}{3},\quad
t = -\frac{\pi}{4},\quad
t = \frac{17\pi}{6},\quad
t = -\frac{7\pi}{2},
\]
after reducing each to a value in $\intoc{-\pi}{\pi}$ modulo $2\pi$.
\end{exercise}

\begin{exercise}[$\star$]\label{exo:g11:trigo:3}
Using the table and the \omterm{prop:g11:trigo:associated}{associated angles}, give the exact values of
\[
\cos\frac{3\pi}{4}, \quad
\sin\frac{5\pi}{6}, \quad
\cos\left(-\frac{\pi}{3}\right), \quad
\sin\frac{4\pi}{3}, \quad
\cos\frac{11\pi}{6} .
\]
\end{exercise}

\begin{exercise}[$\star$]\label{exo:g11:trigo:4}
Given $\cos t = \frac{5}{13}$ with $t \in \intoo{-\frac\pi2}{0}$, compute
$\sin t$ exactly.
\end{exercise}

\begin{exercise}[$\star\star$]\label{exo:g11:trigo:5}
Simplify, for any real $x$:
\[
A = \cos(\pi + x) + \cos(-x) + \sin\left(\frac\pi2 - x\right) - \cos x,
\qquad
B = \sin(\pi - x) + \sin(\pi + x) + \sin(-x) .
\]
\end{exercise}

\begin{exercise}[$\star\star$]\label{exo:g11:trigo:6}
Solve on $\intoc{-\pi}{\pi}$:
\[
\cos t = \frac{\sqrt3}{2}, \qquad
\sin t = \frac12, \qquad
\cos t = -1 .
\]
\end{exercise}

\begin{exercise}[$\star\star$]\label{exo:g11:trigo:7}
Solve on $\intco{0}{2\pi}$:
\[
\sin t = -\frac{\sqrt3}{2}, \qquad
\cos t = 0, \qquad
2\sin t + 1 = 0 .
\]
\end{exercise}

\begin{exercise}[$\star\star$]\label{exo:g11:trigo:8}
Solve on $\R$ the \omterm{def:g10:algebra:equation}{equation} $\cos 2t = \cos t$. (Use
\cref{thm:g11:trigo:equations} with $a = t$, and discuss the two families
of solutions.)
\end{exercise}

\begin{exercise}[$\star\star$]\label{exo:g11:trigo:9}
Show that for all real $x$,
\[
\bigl(\cos x + \sin x\bigr)^2 + \bigl(\cos x - \sin x\bigr)^2 = 2 .
\]
\end{exercise}

\begin{exercise}[$\star\star$]\label{exo:g11:trigo:10}
A wheel of radius $30$ cm rolls without slipping. Through what angle (in
\omterm{def:g11:trigo:radian}{radians}, then in degrees) does it turn when the bike advances by $1.5$ m?
How far does the bike advance during one full turn of the wheel?
\end{exercise}

\begin{exercise}[$\star\star\star$]\label{exo:g11:trigo:11}
Solve on $\intoc{-\pi}{\pi}$ the inequality
$\cos t \leq \frac12$. (Sketch the unit circle, mark the region where the
\omterm{def:g10:coordgeom:system}{abscissa} is at most $\frac12$, and read off the \omterm{def:g10:numbers:interval}{interval} of values of
$t$.)
\end{exercise}
